%!TeX root=../emmatop.tex
\chapter[Chapter \thechapter]{}
\lettrine[lines=4,lraise=0.3]{T}{his} letter must make its way to Emma's feelings. She was obliged, in spite of her previous determination to the contrary, to do it all the justice that Mrs Weston foretold. As soon as she came to her own name, it was irresistible; every line relating to herself was interesting, and almost every line agreeable; and when this charm ceased, the subject could still maintain itself, by the natural return of her former regard for the writer, and the very strong attraction which any picture of love must have for her at that moment. She never stopt till she had gone through the whole; and though it was impossible not to feel that he had been wrong, yet he had been less wrong than she had supposed—and he had suffered, and was very sorry—and he was so grateful to Mrs Weston, and so much in love with Miss Fairfax, and she was so happy herself, that there was no being severe; and could he have entered the room, she must have shaken hands with him as heartily as ever.

She thought so well of the letter, that when Mr Knightley came again, she desired him to read it. She was sure of Mrs Weston's wishing it to be communicated; especially to one, who, like Mr Knightley, had seen so much to blame in his conduct.

»I shall be very glad to look it over,« said he; »but it seems long. I will take it home with me at night.«

But that would not do. Mr Weston was to call in the evening, and she must return it by him.

»I would rather be talking to you,« he replied; »but as it seems a matter of justice, it shall be done.«

He began—stopping, however, almost directly to say, »Had I been offered the sight of one of this gentleman's letters to his mother-in-law a few months ago, Emma, it would not have been taken with such indifference.«

He proceeded a little farther, reading to himself; and then, with a smile, observed, »Humph! a fine complimentary opening: But it is his way. One man's style must not be the rule of another's. We will not be severe.«

»It will be natural for me,« he added shortly afterwards, »to speak my opinion aloud as I read. By doing it, I shall feel that I am near you. It will not be so great a loss of time: but if you dislike it\longdash«

»Not at all. I should wish it.«

Mr Knightley returned to his reading with greater alacrity.

»He trifles here,« said he, »as to the temptation. He knows he is wrong, and has nothing rational to urge.—Bad.—He ought not to have formed the engagement.—»His father«s disposition:'—he is unjust, however, to his father. Mr Weston's sanguine temper was a blessing on all his upright and honourable exertions; but Mr Weston earned every present comfort before he endeavoured to gain it.—Very true; he did not come till Miss Fairfax was here.«

»And I have not forgotten,« said Emma, »how sure you were that he might have come sooner if he would. You pass it over very handsomely—but you were perfectly right.«

»I was not quite impartial in my judgment, Emma:—but yet, I think—had you not been in the case—I should still have distrusted him.«

When he came to Miss Woodhouse, he was obliged to read the whole of it aloud—all that related to her, with a smile; a look; a shake of the head; a word or two of assent, or disapprobation; or merely of love, as the subject required; concluding, however, seriously, and, after steady reflection, thus—

»Very bad—though it might have been worse.—Playing a most dangerous game. Too much indebted to the event for his acquittal.—No judge of his own manners by you.—Always deceived in fact by his own wishes, and regardless of little besides his own convenience.—Fancying you to have fathomed his secret. Natural enough!—his own mind full of intrigue, that he should suspect it in others.—Mystery; Finesse—how they pervert the understanding! My Emma, does not every thing serve to prove more and more the beauty of truth and sincerity in all our dealings with each other?«

Emma agreed to it, and with a blush of sensibility on Harriet's account, which she could not give any sincere explanation of.

»You had better go on,« said she.

He did so, but very soon stopt again to say, »the pianoforte! Ah! That was the act of a very, very young man, one too young to consider whether the inconvenience of it might not very much exceed the pleasure. A boyish scheme, indeed!—I cannot comprehend a man's wishing to give a woman any proof of affection which he knows she would rather dispense with; and he did know that she would have prevented the instrument's coming if she could.«

After this, he made some progress without any pause. Frank Churchill's confession of having behaved shamefully was the first thing to call for more than a word in passing.

»I perfectly agree with you, sir,«—was then his remark. »You did behave very shamefully. You never wrote a truer line.« And having gone through what immediately followed of the basis of their disagreement, and his persisting to act in direct opposition to Jane Fairfax's sense of right, he made a fuller pause to say, »This is very bad.—He had induced her to place herself, for his sake, in a situation of extreme difficulty and uneasiness, and it should have been his first object to prevent her from suffering unnecessarily.—She must have had much more to contend with, in carrying on the correspondence, than he could. He should have respected even unreasonable scruples, had there been such; but hers were all reasonable. We must look to her one fault, and remember that she had done a wrong thing in consenting to the engagement, to bear that she should have been in such a state of punishment.«

Emma knew that he was now getting to the Box Hill party, and grew uncomfortable. Her own behaviour had been so very improper! She was deeply ashamed, and a little afraid of his next look. It was all read, however, steadily, attentively, and without the smallest remark; and, excepting one momentary glance at her, instantly withdrawn, in the fear of giving pain—no remembrance of Box Hill seemed to exist.

»There is no saying much for the delicacy of our good friends, the Eltons,« was his next observation.—»His feelings are natural.—What! actually resolve to break with him entirely!—She felt the engagement to be a source of repentance and misery to each—she dissolved it.—What a view this gives of her sense of his behaviour!—Well, he must be a most extraordinary\longdash«

»Nay, nay, read on.—You will find how very much he suffers.«

»I hope he does,« replied Mr Knightley coolly, and resuming the letter. »»Smallridge!«—What does this mean? What is all this?«

»She had engaged to go as governess to Mrs Smallridge's children—a dear friend of Mrs Elton's—a neighbour of Maple Grove; and, by the bye, I wonder how Mrs Elton bears the disappointment?«

»Say nothing, my dear Emma, while you oblige me to read—not even of Mrs Elton. Only one page more. I shall soon have done. What a letter the man writes!«

»I wish you would read it with a kinder spirit towards him.«

»Well, there is feeling here.—He does seem to have suffered in finding her ill.—Certainly, I can have no doubt of his being fond of her. »Dearer, much dearer than ever.« I hope he may long continue to feel all the value of such a reconciliation.—He is a very liberal thanker, with his thousands and tens of thousands.—»Happier than I deserve.« Come, he knows himself there. »Miss Woodhouse calls me the child of good fortune.«—Those were Miss Woodhouse's words, were they?— And a fine ending—and there is the letter. The child of good fortune! That was your name for him, was it?«

»You do not appear so well satisfied with his letter as I am; but still you must, at least I hope you must, think the better of him for it. I hope it does him some service with you.«

»Yes, certainly it does. He has had great faults, faults of inconsideration and thoughtlessness; and I am very much of his opinion in thinking him likely to be happier than he deserves: but still as he is, beyond a doubt, really attached to Miss Fairfax, and will soon, it may be hoped, have the advantage of being constantly with her, I am very ready to believe his character will improve, and acquire from hers the steadiness and delicacy of principle that it wants. And now, let me talk to you of something else. I have another person's interest at present so much at heart, that I cannot think any longer about Frank Churchill. Ever since I left you this morning, Emma, my mind has been hard at work on one subject.«

The subject followed; it was in plain, unaffected, gentlemanlike English, such as Mr Knightley used even to the woman he was in love with, how to be able to ask her to marry him, without attacking the happiness of her father. Emma's answer was ready at the first word. »While her dear father lived, any change of condition must be impossible for her. She could never quit him.« Part only of this answer, however, was admitted. The impossibility of her quitting her father, Mr Knightley felt as strongly as herself; but the inadmissibility of any other change, he could not agree to. He had been thinking it over most deeply, most intently; he had at first hoped to induce Mr Woodhouse to remove with her to Donwell; he had wanted to believe it feasible, but his knowledge of Mr Woodhouse would not suffer him to deceive himself long; and now he confessed his persuasion, that such a transplantation would be a risk of her father's comfort, perhaps even of his life, which must not be hazarded. Mr Woodhouse taken from Hartfield!—No, he felt that it ought not to be attempted. But the plan which had arisen on the sacrifice of this, he trusted his dearest Emma would not find in any respect objectionable; it was, that he should be received at Hartfield; that so long as her father's happiness—in other words, his life—required Hartfield to continue her home, it should be his likewise.

\begin{figure}[tbph]
\centering
\includegraphics[width=\linewidth]{51walking}
\caption{Walking away from William Larkins}
\end{figure}

Of their all removing to Donwell, Emma had already had her own passing thoughts. Like him, she had tried the scheme and rejected it; but such an alternative as this had not occurred to her. She was sensible of all the affection it evinced. She felt that, in quitting Donwell, he must be sacrificing a great deal of independence of hours and habits; that in living constantly with her father, and in no house of his own, there would be much, very much, to be borne with. She promised to think of it, and advised him to think of it more; but he was fully convinced, that no reflection could alter his wishes or his opinion on the subject. He had given it, he could assure her, very long and calm consideration; he had been walking away from William Larkins the whole morning, to have his thoughts to himself.

»Ah! there is one difficulty unprovided for,« cried Emma. »I am sure William Larkins will not like it. You must get his consent before you ask mine.«

She promised, however, to think of it; and pretty nearly promised, moreover, to think of it, with the intention of finding it a very good scheme.

It is remarkable, that Emma, in the many, very many, points of view in which she was now beginning to consider Donwell Abbey, was never struck with any sense of injury to her nephew Henry, whose rights as heir-expectant had formerly been so tenaciously regarded. Think she must of the possible difference to the poor little boy; and yet she only gave herself a saucy conscious smile about it, and found amusement in detecting the real cause of that violent dislike of Mr Knightley's marrying Jane Fairfax, or any body else, which at the time she had wholly imputed to the amiable solicitude of the sister and the aunt.

This proposal of his, this plan of marrying and continuing at Hartfield—the more she contemplated it, the more pleasing it became. His evils seemed to lessen, her own advantages to increase, their mutual good to outweigh every drawback. Such a companion for herself in the periods of anxiety and cheerlessness before her!—Such a partner in all those duties and cares to which time must be giving increase of melancholy!

She would have been too happy but for poor Harriet; but every blessing of her own seemed to involve and advance the sufferings of her friend, who must now be even excluded from Hartfield. The delightful family party which Emma was securing for herself, poor Harriet must, in mere charitable caution, be kept at a distance from. She would be a loser in every way. Emma could not deplore her future absence as any deduction from her own enjoyment. In such a party, Harriet would be rather a dead weight than otherwise; but for the poor girl herself, it seemed a peculiarly cruel necessity that was to be placing her in such a state of unmerited punishment.

In time, of course, Mr Knightley would be forgotten, that is, supplanted; but this could not be expected to happen very early. Mr Knightley himself would be doing nothing to assist the cure;—not like Mr Elton. Mr Knightley, always so kind, so feeling, so truly considerate for every body, would never deserve to be less worshipped than now; and it really was too much to hope even of Harriet, that she could be in love with more than three men in one year.